\documentclass[10pt,a4paper]{amsart}
\usepackage[margin=19mm]{geometry}
\usepackage{amsmath,amssymb,mathtools}
\usepackage[scr=rsfs,cal=euler]{mathalfa}
\usepackage{xfrac}
\usepackage[unicode,hidelinks,bookmarks=false]{hyperref}
\newcommand{\ie}{i.e.\ }
\newcommand{\cf}{cf.\ }
\newcommand{\resp}{respectively\ }
\newcommand{\Subsection}{\subsection}
\providecommand{\nobreakdash}{\nobreak\hbox{-}}
\setlength{\emergencystretch}{2em}
\multlinegap0pt
\usepackage{amssymb}
\usepackage{mathtools}
\usepackage[shortlabels]{enumitem}
\usepackage{xcolor}
\newtheorem{lemma}{Lemma}
\newcommand{\BS}{\mathrm{BS}}
\title{Accessible CAT$(-1)$ groups of critical exponent less than one}
\author{Yong Hou}
\date{Source: arXiv:2609.00557v1 (CC BY 4.0)}
\begin{document}
\maketitle

\noindent\emph{Source and licence.} Accessible CAT$(-1)$ groups of critical exponent less than one. Version arXiv:2609.00557v1 (CC BY 4.0), distributed by arXiv under the Creative Commons Attribution 4.0 International licence, http://creativecommons.org/licenses/by/4.0/, as recorded in the arXiv licence metadata for this exact version and shown on the abstract page https://arxiv.org/abs/2609.00557v1. Full licence notice: \texttt{licenses/arxiv-2609.00557-CC-BY-4.0.txt}.\par\smallskip

\noindent\textit{Editorial scope.} Counted unit: the article's lemma lem:bass-serre-normal-form (source0.tex lines 2018--2033). The article's complete original proof of it is delivered with it (source0.tex lines 2035--2102, one \texttt{proof} environment of 68 lines). No mathematics is written, repaired or re-derived: the statement and the proof are the article's own lines, in the article's own notation, and no retained line is substituted or re-flowed.  Retained uncounted as the source-local context the proof needs: lines 1952--1956; lines 1965--1967.  Nothing was omitted from any retained span, so the package carries no omission record.  No retained line contains a citation, so the wrapper carries no bibliography and no bibliography entry.  No retained line contains a figure environment, an image-inclusion command or drawing code, so no image is delivered and the package declares no asset dependency.  Editorial formatting only, in addition: an AMS \texttt{amsart} preamble that restates the article's own declarations and macros used by the retained lines, namely the environments \texttt{lemma} on separate counters, together with the standard packages those lines need.  The retained environments print in this document as Lemma 1 in order of appearance, whereas the article prints the same items with the numbering of its own full document.  The complete native source of the article as distributed in the arXiv e-print for this exact version is bundled unchanged as \texttt{sources/209039/source0.tex}.
\begin{equation}\label{eq:reverse-edge-lifts}
 \widetilde{\bar f}=\tau_f^{-1}\overline{\widetilde f},
 \qquad
 \tau_{\bar f}=\tau_f^{-1}.
\end{equation}

\begin{equation}\label{eq:reverse-edge-stabilizer}
 G_{\bar f}=\tau_f^{-1}G_f\tau_f.
\end{equation}

\begin{lemma}
\label{lem:bass-serre-normal-form}
Every $g\in H$ has a unique representation
\begin{equation}\label{eq:bass-serre-normal-form}
 g=b(\omega)h,
 \qquad h\in G_{v_*},
\end{equation}
where $\omega$ is a reduced coset-labelled loop based at $v_*$.  Consequently
\begin{equation}\label{eq:bass-serre-length}
 L^{\BS}_x(g)=\sum_{j=1}^nL_{f_j,x}(c_j)+d(x,hx)
\end{equation}
defines a normal form length and
\begin{equation}\label{eq:bass-serre-subadditive-upper}
 d(x,gx)\le L^{\BS}_x(g).
\end{equation}
\end{lemma}

\begin{proof}
Write $x_*=\widetilde v_*$.  For a relative coset
$c_j=r_jG_{f_j}$ set
\[
 a_j=r_jh_{f_j}(c_j)\in G_{o(f_j)},
 \qquad
 p_0=1,
 \qquad
 p_j=b_{f_1}(c_1)\cdots b_{f_j}(c_j)
      =p_{j-1}a_j\tau_{f_j}.
\]
The representative $a_j$ belongs to the coset $c_j$, and multiplication by an element of
$G_{f_j}$ does not change the actual lift of $f_j$.  Hence, whenever the current vertex is
$p_{j-1}\widetilde{o(f_j)}$, the coset $c_j$ determines the actual oriented edge
\begin{equation}\label{eq:actual-bass-serre-edge}
 E_j=p_{j-1}a_j\widetilde f_j.
\end{equation}
Its terminal vertex is
\begin{equation}\label{eq:actual-bass-serre-terminal}
 t(E_j)=p_{j-1}a_j\tau_{f_j}\widetilde{t(f_j)}
       =p_j\widetilde{t(f_j)}.
\end{equation}
These two identities are the induction that relates the coordinate word to the path in $T$.

We next verify the non-backtracking rule exactly.  Suppose
$f_{j+1}=\bar f_j$.  By \eqref{eq:reverse-edge-lifts}, the next actual edge is
\[
 p_j a_{j+1}\widetilde{\bar f_j}
 =p_{j-1}a_j\tau_{f_j}a_{j+1}\tau_{f_j}^{-1}
   \overline{\widetilde f_j}.
\]
This equals the reverse
$\overline{E_j}=p_{j-1}a_j\overline{\widetilde f_j}$ if and only if
$\tau_{f_j}a_{j+1}\tau_{f_j}^{-1}\in G_{f_j}$, or equivalently, by
\eqref{eq:reverse-edge-stabilizer}, if and only if
$a_{j+1}\in G_{\bar f_j}$.  This is precisely the identity relative coset in
$G_{o(\bar f_j)}/G_{\bar f_j}$.  Thus omitting that coset is exactly the condition that the
actual tree path have no immediate reversal.

Now let $g\in H$.  The geodesic edge path from $x_*$ to $gx_*$ projects to an edge loop
$f_1\cdots f_n$ based at $v_*$.  Inductively, after the first $j-1$ edges have been encoded,
the $j$th actual outgoing edge is uniquely of the form
$p_{j-1}r\widetilde f_j$ for one left coset
$rG_{f_j}\in G_{o(f_j)}/G_{f_j}$.  Choose this coset as $c_j$ and use its fixed minimising
representative $a_j$. Equations \eqref{eq:actual-bass-serre-edge} and
\eqref{eq:actual-bass-serre-terminal} show that the coordinate word follows the same actual
edge and reaches the same terminal vertex.  The geodesic has no immediate reversal, so the
coset-labelled loop is reduced by the preceding calculation.  At the endpoint,
$p_nx_*=gx_*$, and consequently
$h=p_n^{-1}g$ belongs to $G_{v_*}$.  This proves existence of
\eqref{eq:bass-serre-normal-form}.

Conversely, equations \eqref{eq:actual-bass-serre-edge} and
\eqref{eq:actual-bass-serre-terminal} associate to every reduced coset-labelled loop a path
in $T$ from $x_*$ to $p_nx_*$.  The calculation above shows that no two consecutive edges
are reverses, so this path is reduced and hence geodesic.  Suppose
$p_nh=p'_mh'$ are two representations of the same element.  Since $h,h'$ fix $x_*$, the two
reduced paths have the same endpoints.  A tree has a unique geodesic between two vertices,
so the actual edge paths agree.  Their projections give $n=m$ and
$f_j=f'_j$ for every $j$.  Starting with $p_0=p'_0=1$, equality of the $j$th actual outgoing
edges gives equality of the corresponding left cosets, hence $c_j=c'_j$; because the
representative $b_{f_j}(c_j)$ was fixed once and for all, it then gives $p_j=p'_j$.
Induction recovers all cosets and finally $h=h'$.  This proves uniqueness.

Finally, subadditivity of orbit length and the minimising property
$d(x,b_{f_j}(c_j)x)=L_{f_j,x}(c_j)$ give
$d(x,gx)\le L_x^{\BS}(g)$.
\end{proof}

\end{document}
